\section{Extension across a mixed boundary}\label{sec:boundary}

We next extend a box across a smooth boundary point whose complex tangent
involves both source directions.  The K\"ahler form controls the areas of
shrinking slices.  An annulus of small area then supplies the shell of a
Hartogs figure large enough to cross the boundary.  We first isolate the
elementary compactness argument that converts small area into a small
image. Its proof is a local form of the derivative normalization used in
Brody's reparametrization argument \cite{Brody78}, with vanishing area
providing the contradiction.

\begin{lemma}[Small area]\label{lem:small-area}
Let $M$ be a compact Hermitian manifold, with associated positive
$(1,1)$-form $\omega_M$, and let $U\subset\C$ be a domain.  Suppose that
$h_j:U\to M$ are holomorphic and
\[
  \int_U h_j^*\omega_M\longrightarrow 0.
\]
Then $\norm{\dd h_j}$ tends to zero uniformly on every compact subset of $U$.
Here the source has its Euclidean metric.  Closedness of $\omega_M$ is not
required.
\end{lemma}

\begin{proof}
It suffices to work on a disk $D$ with $\overline D\subset U$.  Put
$d(z)=\operatorname{dist}(z,\partial D)$ and
\[
  M_j=\max_{\overline D}d(z)\norm{\dd h_j(z)}.
\]
If $M_j\not\to0$, pass to a subsequence with $M_j\ge\varepsilon>0$.
Choose a maximizing point $z_j\in D$, and set
$a_j=\norm{\dd h_j(z_j)}$ and $d_j=d(z_j)$, so $a_jd_j=M_j$.
The maps
\[
  k_j(\zeta)=h_j(z_j+\zeta/a_j)
\]
are defined for $|\zeta|<\varepsilon/4$.  On this disk,
$d(z_j+\zeta/a_j)\ge d_j/2$, and maximality gives
\[
  \norm{\dd k_j(\zeta)}\le2,
  \qquad \norm{\dd k_j(0)}=1.
\]
After taking a subsequence, $k_j(0)$ converges to a point of $M$.
The derivative bound puts the images of a fixed smaller disk in a
relatively compact coordinate neighborhood of that point.  Montel's
theorem and the Cauchy estimates give a further subsequence converging
there together with its derivatives.  Its limit has derivative norm $1$
at the origin and therefore positive area on a smaller disk.  This
contradicts the vanishing areas of $k_j$, which are bounded by the areas
of $h_j$ and are unchanged by the affine reparametrization.  Thus
$M_j\to0$.  On each smaller disk $d$ is bounded below, so the derivatives
tend uniformly to zero.  A finite covering proves the assertion on any
compact subset of $U$.
\end{proof}

\begin{proposition}[Mixed-boundary extension]\label{prop:mixed}
Let $B\subset\C^2_{s,t}$ be a convex open set, and let $f:B\to X$ be a
holomorphic box.  Suppose that $p\in\partial B$ has a $C^2$ defining
function with nonzero differential.  If the complex tangent line to
$\partial B$ at $p$ is generated by a vector with both components
nonzero, then $f$ extends as a holomorphic box to a neighborhood of $p$,
agreeing with $f$ on its intersection with $B$.
\end{proposition}

\begin{proof}
Use the fixed K\"ahler form $\omega$ and its Hermitian norm, and recall
the bounded projections $P_i$ from \eqref{eq:projection-bounds}.

Write $p=(s_*,t_*)$.  Choose complex-affine coordinates $(u,v)$ with
\begin{equation}\label{eq:boundary-affine}
  (s,t)=(s_*,t_*)+(a_1,a_2)u+(b_1,b_2)v,
  \qquad b_1b_2\ne0,
\end{equation}
where the $v$-direction is complex tangent and the positive real
$u$-direction points outward.  All constants below may depend on this
boundary point and these coordinates.

\paragraph{Shrinking slices and bounded area.}
Take a local defining function $\rho$ with $B=\{\rho<0\}$ near $p$.
For real $x$ and small $v$, Taylor's theorem gives
\[
  \rho(x,v)\le c x+M(x^2+|v|^2),\qquad c>0,
\]
because both real derivatives in the complex tangent direction vanish.
Choose $k>0$ sufficiently small.  For all $x<0$ sufficiently close to
zero, the closed disks
\[
  D_x=\{u=x,\ |v|\le r(x)\},\qquad r(x)=k\sqrt{-x},
\]
are compactly contained in $B$.

Let $e_u=\partial_{\operatorname{Re}u}$ and
$e_v=\partial_{\operatorname{Re}v}$.  The box tangencies and complex
linearity give the exact identity
\[
  \dd f(e_u)=\sum_{i=1}^2\frac{a_i}{b_i}P_i\dd f(e_v).
\]
Consequently
\begin{equation}\label{eq:boundary-speed}
  \norm{\dd f(e_u)}\le C\norm{\dd f(e_v)},\qquad
  C=\sum_{i=1}^2\left|\frac{a_i}{b_i}\right|
       \sup_X\norm{P_i}.
\end{equation}
Holomorphicity makes $\norm{\dd f(e_v)}$ equal to the speed in every unit
real direction of the $v$-plane, including the radial direction.

Set $\Omega=f^*\omega$ and $A(x)=\int_{D_x}\Omega$.  Closedness of
$\Omega$ and Stokes' theorem, applied to the moving disks, yield
\[
 A'(x)=\int_0^{2\pi}
   \Omega\bigl(e_u+r'(x)\partial_r,\partial_\theta\bigr)
       \big|_{u=x,\ v=r(x)e^{i\theta}}\,\dd\theta.
\]
Indeed the boundary velocity is $e_u+r'(x)\partial_r$, and the circles
have their counterclockwise orientation.  On a boundary circle put
$q=\norm{\dd f(e_v)}$.  Since $\partial_\theta=rJ\partial_r$,
\[
  \Omega(\partial_r,\partial_\theta)=rq^2,
  \qquad
  |\Omega(e_u,\partial_\theta)|\le Crq^2.
\]
Thus
\[
  A'(x)\le\int_0^{2\pi}r(x)\bigl(C+r'(x)\bigr)q^2\,\dd\theta.
\]
Here $r'(x)=-k/(2\sqrt{-x})$.  Choose $x_0<0$ sufficiently close to
zero that $r'(x)\le-C$ for $x_0\le x<0$.  We obtain the uniform bound
\begin{equation}\label{eq:boundary-area}
  0\le A(x)\le A(x_0)<\infty\qquad(x_0\le x<0).
\end{equation}
No product property of the K\"ahler metric has been used.

\paragraph{A small-image annulus.}
Let $x_j=-\delta_j$ with $\delta_j\downarrow0$.  Among the scales
$\lambda=5^\ell\delta_j^{3/4}$, with integer $\ell\ge0$, in
$[\delta_j^{3/4},\delta_j^{5/8}]$, the annuli
\[
  \lambda/2<|v|<2\lambda
\]
are pairwise disjoint, and their number tends to infinity.  For large
$j$ they lie in $D_{x_j}$.  By \eqref{eq:boundary-area}, one of them,
with scale $\lambda_j$, has area tending to zero.  These choices satisfy
\begin{equation}\label{eq:boundary-scales}
  \frac{\delta_j}{\lambda_j}\le\delta_j^{1/4}\longrightarrow0,
  \qquad
  \frac{\lambda_j}{\sqrt{\delta_j}}
       \le\delta_j^{1/8}\longrightarrow0.
\end{equation}
Apply Lemma~\ref{lem:small-area} to
$h_j(\zeta)=f(x_j,\lambda_j\zeta)$ on $1/2<|\zeta|<2$.
The derivatives tend uniformly to zero on a neighborhood of
$3/4\le|\zeta|\le3/2$.  This compact annulus has bounded intrinsic
diameter, so the image of
\[
  u=x_j,\qquad 3\lambda_j/4<|v|<3\lambda_j/2
\]
has diameter tending to zero.  A finite cover of $X$ by product charts
has a Lebesgue number.  Hence, for large $j$, this entire image lies in
one product chart $z:V\to V_1\times V_2$.  Write its coordinate functions
as $(g_1(v),g_2(v))$.

Figure~\ref{fig:boundary-scales} separates the two scale comparisons.
\begin{figure}[tbp]
\centering
\setlength{\unitlength}{1pt}
\begin{picture}(430,145)
  \put(100,125){\makebox(0,0){\small normal real slice ($v=0$)}}
  \put(15,75){\vector(1,0){180}}
  {\color{blue!65!black}\linethickness{2pt}\put(35,75){\line(1,0){130}}}\ignorespaces
  \put(100,71){\line(0,1){8}}
  \put(116,68){\line(0,1){14}}
  \put(100,75){\circle*{3}}
  \put(116,75){\circle*{3}}
  \put(91,58){\makebox(0,0){\small $x_j=-\delta_j$}}
  \put(127,91){\makebox(0,0){\small $p: u=0$}}
  \put(35,101){\line(1,0){130}}
  \put(35,98){\line(0,1){6}}
  \put(165,98){\line(0,1){6}}
  \put(100,111){\makebox(0,0){\small shell width $2\eta\lambda_j$}}
  \put(193,63){\makebox(0,0){\small $u$}}
  \put(108,37){\makebox(0,0){\small $\delta_j/\lambda_j\longrightarrow0$}}
  \put(320,125){\makebox(0,0){\small complex tangent slice ($u=x_j$)}}
  \put(320,69){\circle{90}}
  {\color{blue!65!black}\put(320,69){\circle{34}}\put(320,69){\circle{43}}}\ignorespaces
  \put(320,69){\circle*{2}}
  \put(320,69){\vector(1,0){45}}
  \put(376,58){\makebox(0,0)[l]{\small $r(x_j)$}}
  \put(320,69){\vector(-4,3){16}}
  \put(300,90){\makebox(0,0)[r]{\small $\lambda_j$}}
  \put(320,8){\makebox(0,0){\small $\lambda_j/\sqrt{\delta_j}\longrightarrow0$}}
\end{picture}
\caption{The two scale inequalities used to cross the boundary. On the
left, the filled Hartogs domain extends in the $u$ direction farther than
the distance $\delta_j$ from the slice to $p$. On the right, its shell
has radius comparable to $\lambda_j$ and fits well inside the slice disk
of radius $r(x_j)=k\sqrt{\delta_j}$. The annulus has a small image in one
target product chart. These are separate source slices in the affine
coordinates $(u,v)$, drawn schematically and not at a common scale.}
\label{fig:boundary-scales}
\end{figure}


We have found an annulus whose radius is much larger than the remaining
distance $\delta_j$ to the boundary.  We now extend it transversely to
make a Hartogs shell.

\paragraph{The shell and its agreement with the old box.}
Put $w=u-x_j$ and $c_i=a_i/b_i$.  Choose $\eta>0$, independent of $j$,
so that $\eta\max(1,|c_1|,|c_2|)<1/8$.  On
\[
  S_j=\{|w|<\eta\lambda_j,\quad
                \lambda_j<|v|<5\lambda_j/4\}
\]
define
\begin{equation}\label{eq:boundary-shell}
  F_j(w,v)=z^{-1}\bigl(g_1(v+c_1w),g_2(v+c_2w)\bigr).
\end{equation}
Both shifted arguments stay in
$3\lambda_j/4<|v|<3\lambda_j/2$.  The range of $z$ is a product, so
combining the two valid factor coordinates in
\eqref{eq:boundary-shell} needs no further target-chart margin.  Moreover,
\[
  v+c_1w=\frac{s-s_*-a_1x_j}{b_1},\qquad
  v+c_2w=\frac{t-t_*-a_2x_j}{b_2}.
\]
Thus $F_j$ is a box in the original $(s,t)$ directions.

The two maps agree on an open germ at every point of $S_j$ with $w=0$.
To see this, continuity puts the old map $f$ in $V$ on a neighborhood
of that point.  There its separated expression is
$(\varphi_1(s),\varphi_2(t))$.  Restriction to $u=x_j$ gives the local
expressions for $g_i$, and the displayed identities show that
\eqref{eq:boundary-shell} reproduces $f$ on a full neighborhood.  This
is stronger than equality merely on the slice $w=0$.

The closed slice disk $u=x_j$, $|v|\le5\lambda_j/4$ is contained in
$B$.  Hence there is $0<e_j<\eta\lambda_j$ such that the original map
is defined on
\[
  C_j=\{|w|<e_j,\ |v|<5\lambda_j/4\}.
\]
The overlap $C_j\cap S_j$ is a connected product of a disk and an
annulus.  The identity theorem for maps into $X$ therefore extends the
open-germ agreement to this entire overlap, even if the old map elsewhere
leaves $V$.  We have a box on the Hartogs figure $C_j\cup S_j$.

By Proposition~\ref{prop:hartogs}, it extends to a box on
\[
  Q_j=\{|w|<\eta\lambda_j,\ |v|<5\lambda_j/4\}.
\]
Both $Q_j$ and $B$ are convex in the underlying real affine space.
Their intersection is connected and contains $C_j$, so the extension
agrees with $f$ throughout $Q_j\cap B$.  Finally $p$ has coordinates
$(w,v)=(\delta_j,0)$, which belong to $Q_j$ for large $j$ by
\eqref{eq:boundary-scales}.  This gives the required extension across
$p$.
\end{proof}
